\chapter{Presentaci\'on del PFC BIOPET}
\label{cap:presentacion-pfc}

\section{Qu\'e es BIOPET}
\label{sec:que-es-biopet}

BIOPET es un sistema web de gesti\'on veterinaria, desarrollado como
Proyecto Fin de Curso (PFC) de la asignatura \asignatura\ por su equipo
original (Mariscal Cabrera Jaime Josu\'e, Beltr\'an Montiel Fred Adri\'an y
Taipe Mora Zaida Melissa; ver secci\'on~\ref{sec:distincion-equipos}).
Permite el registro y autenticaci\'on de usuarios, y la administraci\'on de
mascotas, citas veterinarias, consultas (historial cl\'inico) y vacunas,
con control de acceso diferenciado por rol y, en varios recursos, por
propiedad del registro.

\section{Prop\'osito y alcance funcional}
\label{sec:proposito-alcance}

El sistema define cuatro roles reales, verificados en el c\'odigo
(\texttt{Rol}, columna \texttt{usuarios.rol}): \texttt{ROLE\_ADMIN},
\texttt{ROLE\_VETERINARIO}, \texttt{ROLE\_AUXILIAR} y \texttt{ROLE\_DUENO}.
Los tres primeros tienen alcance global sobre los recursos gestionados por
la cl\'inica; \texttt{ROLE\_DUENO} est\'a limitado, en la mayor\'ia de los
recursos, a consultar \'unicamente lo asociado a sus propias mascotas.

Los m\'odulos funcionales actualmente implementados y verificados contra el
c\'odigo del backend (\texttt{Backend/src/main/java/com/biopet/controller/})
son:

\begin{itemize}[nosep]
  \item \textbf{Autenticaci\'on} (\texttt{AuthController}): registro,
    login, refresh y logout, mediante JWT entregado por cookies
    \texttt{HttpOnly}/\texttt{Secure}/\texttt{SameSite=Strict}.
  \item \textbf{Usuarios} (\texttt{UsuarioController}): consulta del
    perfil propio (\texttt{/me}) y CRUD administrativo completo, este
    \'ultimo restringido a \texttt{ROLE\_ADMIN}.
  \item \textbf{Mascotas} (\texttt{MascotaController}): CRUD completo con
    control de acceso por rol y por propiedad, m\'as un resumen
    estad\'istico por especie respaldado por una funci\'on almacenada de
    PostgreSQL.
  \item \textbf{Citas} (\texttt{CitaController}): CRUD de citas
    veterinarias, asociadas a una mascota y a un veterinario asignado.
  \item \textbf{Consultas} (\texttt{ConsultaController}): registro de
    diagn\'ostico, tratamiento y observaciones de una atenci\'on m\'edica,
    asociado cronol\'ogicamente al historial de la mascota.
  \item \textbf{Vacunas} (\texttt{VacunaController}): registro de vacunas
    aplicadas por mascota, con consulta filtrable por mascota.
  \item \textbf{Integraci\'on con una API externa}
    (\texttt{ExternalApiController}): consulta de informaci\'on de una
    especie contra un servicio de terceros, con cach\'e Redis
    (patr\'on \emph{cache-aside}).
\end{itemize}

\section{Arquitectura y stack tecnol\'ogico}
\label{sec:arquitectura-stack}

BIOPET sigue una arquitectura de tres capas m\'as cach\'e: un frontend SPA
en Angular, un backend REST en Spring Boot, una base de datos relacional
PostgreSQL y una capa de cach\'e/revocaci\'on de tokens en Redis, todo
orquestado con Docker Compose.

\begin{table}[H]
  \centering
  \caption{Stack tecnol\'ogico de BIOPET, verificado contra el c\'odigo actual.}
  \label{tab:stack-tecnologico}
  \begin{tabular}{@{}ll@{}}
    \toprule
    Componente & Versi\'on / detalle \\
    \midrule
    Backend & Java 21, Spring Boot 3.2.12~\cite{spring-boot-docs} \\
    Seguridad & Spring Security 6, JWT (\texttt{jjwt} 0.12.6, HMAC-SHA256)~\cite{rfc7519} \\
    Persistencia & Spring Data JPA + Hibernate, PostgreSQL 16, Flyway \\
    Cach\'e / revocaci\'on & Redis 7 (Spring Data Redis + Spring Cache) \\
    Documentaci\'on de API & springdoc-openapi 2.5.0 (Swagger UI)~\cite{openapi-spec} \\
    Cobertura & JaCoCo 0.8.12 \\
    Frontend & Angular 17.3, TypeScript 5.4, componentes standalone~\cite{angular-docs} \\
    Manejo uniforme de errores & \texttt{ProblemDetail} (RFC 7807)~\cite{rfc7807} \\
    Orquestaci\'on & Docker Compose, \texttt{Makefile} \\
    \bottomrule
  \end{tabular}
\end{table}

\evidencia{../../diagrams/c4-contenedores/c4-contenedores.png}{Diagrama C4 Nivel 2 (contenedores) de BIOPET.}{fig:c4-contenedores}

El detalle completo de la arquitectura, con relaciones verificadas contra
el c\'odigo, sigue el modelo C4~\cite{c4-model} y est\'a documentado en
\texttt{docs/diagrams/c4-contenedores/} (Nivel 2, contenedores) y
\texttt{docs/diagrams/c4-componentes-backend/} (Nivel 3, componentes del
backend), ambos parte del repositorio heredado del PFC.

\section{Estado del proyecto al momento de esta revisi\'on}
\label{sec:estado-proyecto}

Seg\'un el \texttt{README.md} ra\'iz del repositorio, BIOPET se presenta
como \textbf{versi\'on final acad\'emica} para el alcance definido en esta
actividad de Unidad IV: las funcionalidades listadas en la
secci\'on~\ref{sec:proposito-alcance} est\'an integradas, las pruebas
automatizadas pasan (ver cap\'itulo~\ref{cap:resultados}), y existe
documentaci\'on t\'ecnica y evidencia reproducible respaldando cada
componente. El propio README es expl\'icito en que esto \textbf{no
equivale a certificar el sistema como listo para un entorno productivo
real}: usa un certificado TLS autofirmado y credenciales de desarrollo, y
no fue evaluado en un entorno equivalente a producci\'on.
